By default, \edps\ associates raw calibrations to the reduction
process. It is also possible to use pre-processed calibrations
(a.k.a. master calibrations) if available, in order to speed up the
reduction.  The preference can be specified in the Raw Data tab,
before creating the datasets.



Possible values of the Calibration Preferences are:

\begin{itemize}

 \item {\bf raw\_per\_quality\_level}: At equal quality of reduction,
   association of raw calibrations is preferred. This is the default.

\item {\bf master\_per\_quality\_level}: At equal quality of
  reduction, association of master calibrations is preferred.

\item {\bf raw}: Association of raw calibration is preferred, despite
  the quality of results.
  
\item {\bf master}: Association of master calibration is preferred,
  despite the quality of results.
\end{itemize}

When master calibrations are used, the reduction step needed to
process raw calibrations are not executed. The reduction then moves
directly to the process of scientific exposures.

For example, if reduction speed for a quick check is preferred over a
high quality reduction, one can select "master".  In this case, old
master calibrations are associated even if there are raw calibrations
closer in time (and therefore more likely to ensure better quality
products).

The quality level that the selected calibrations deliver is indicated
close to each dataset in the `Raw input` tab, under the colum
`CalibLevel`. CalibLevel=0 indicates that calibrations that follow the
rules of the instrument calibration plans have been selected. The
higher the number, the poorer the quality of the products.

